% GENERATED FILE. Edit canonical lessons, then run npm run generate:books.
% canonical-source-hash: fnv1a64:873301090dfd37b5
% canonical-lessons: HI-C72-money, HI-C72-rupee, HI-C72-price, HI-C72-expensive, HI-C72-buy

\chapter{What It Costs}
\label{ch:hi-what-it-costs}
\hlchaptermodality{79}

\begin{chapteropening}
\textbf{By the end of this chapter:} \emph{I can ask what something costs, say that it is expensive, and say that I am buying it.}

Complete HI-C72-buy and say all five words in order, then ask what a thing costs and call it expensive.
\end{chapteropening}

\section[paisā]{\hi{पैसा} (paisā) — money}

\label{lesson:HI-C72-money}



\begin{quote}
Two words back: what does \emph{bhī} mean, and what is \textbf{\hi{बहुत}}? The next five words are what you need at a counter.
\end{quote}

\subsection*{You'll want to know: \hi{पैसा}}
\textbf{\hi{पैसा}} (\emph{paisā}) — \textquotedblleft{}money\textquotedblright{}.

It began as a fraction. Sanskrit \hi{पाद} (\emph{pāda}) is a quarter and \hi{अंश} (\emph{aṁśa}) is a part, and the coin they named was a small share of a larger one. For a long while a \hi{पैसा} was one hundredth of a rupee.

Then the small coin ate the whole idea. Today \hi{पैसा} as one coin is nearly gone, and \hi{पैसा} as \textbf{money in general} is what people say: \emph{paisā nahīṁ hai} is \textquotedblleft{}there is no money\textquotedblright{}, about a wallet or about a country.

The \hi{ै} on the \hi{प} is the \emph{ai} matra, two strokes above the headline, and it is the same one inside \textbf{\hi{कैसे}}.

\begin{grammarlens}[title={What you've built}]
The first of five words for buying something.
\end{grammarlens}

\subsection*{Your turn}
\begin{itemize}
\raggedright
  \item \emph{Recall:} say \emph{bhī}, then read \textbf{\hi{बहुत}} and say what it means
  \item \emph{Hear:} \emph{paisā}
  \item \emph{Say it:} \emph{paisā}
  \item \emph{Read it:} \hi{पैसा}, and find the \hi{ै} it shares with \hi{कैसे}
\end{itemize}

\begin{tcolorbox}[breakable,colback=teal!4,colframe=teal!35!black,title={Before you move on}]
What does \hi{पैसा} mean? (\textquotedblleft{}money\textquotedblright{}.) Say it once more.
\end{tcolorbox}
\section[rupayā]{\hi{रुपया} (rupayā) — a rupee}

\label{lesson:HI-C72-rupee}



\begin{quote}
Say \emph{paisā}, then write \textbf{\hi{पैसा}}: \hi{प} with \hi{ै} above it, then \hi{स}, then \hi{ा}.
\end{quote}

\subsection*{You'll want to know: \hi{रुपया}}
\textbf{\hi{रुपया}} (\emph{rupayā}) — \textquotedblleft{}a rupee\textquotedblright{}.

Sanskrit \hi{रूप्य} (\emph{rūpya}) is \textquotedblleft{}wrought silver, a stamped coin\textquotedblright{}, and it comes straight from \textbf{\hi{रूप}} (\emph{rūpa}), \textquotedblleft{}form, shape\textquotedblright{}. A rupee is a piece of metal that has been \textbf{given a shape} — the name says nothing about value and everything about the stamping.

That is worth holding beside \hi{पैसा}. \hi{पैसा} was named for being a fraction; \hi{रुपया} was named for being formed. One word looks at how small it is, the other at how it was made.

The \hi{ु} under the \hi{र} is the short \emph{u} matra, hanging below the letter rather than above it.

A price is said in \hi{रुपये}: \textbf{\hi{दस} \hi{रुपये}}, and you have counted to ten already.

\begin{grammarlens}[title={What you've built}]
Money in general, and the unit it is counted in.
\end{grammarlens}

\subsection*{Your turn}
\begin{itemize}
\raggedright
  \item \emph{Recall:} read \textbf{\hi{बहुत}}, then say it without looking
  \item \emph{Say it:} \emph{rupayā}
  \item \emph{Say it:} \emph{das rupaye}
  \item \emph{Write it:} \hi{पैसा} once more, then \hi{रुपया} beside the printed model
\end{itemize}

\begin{tcolorbox}[breakable,colback=teal!4,colframe=teal!35!black,title={Before you move on}]
What does \hi{रुपया} mean? (\textquotedblleft{}a rupee\textquotedblright{}.) And \hi{पैसा}? (\textquotedblleft{}money\textquotedblright{}.)
\end{tcolorbox}
\section[dām]{\hi{दाम} (dām) — a price}

\label{lesson:HI-C72-price}



\begin{quote}
Say the unit a price is counted in, and then say ten of them.
\end{quote}

\subsection*{You'll want to know: \hi{दाम}}
\textbf{\hi{दाम}} (\emph{dām}) — \textquotedblleft{}a price\textquotedblright{}.

This one has travelled further than any word in the chapter. The Greek coin was the \emph{drachmē}; traders carried the name east; Sanskrit wrote it \hi{द्रम्म} (\emph{dramma}); Hindi wore it down to \hi{दाम}. A Greek weight of silver ended up as the ordinary Hindi word for what a thing costs.

Now put it with a word you have had since the second chapter. \textbf{\hi{क्या}} is \textquotedblleft{}what\textquotedblright{}, and

\begin{quote}
\textbf{\hi{दाम} \hi{क्या} \hi{है}?} — \emph{dām kyā hai?} — \textquotedblleft{}what is the price?\textquotedblright{}
\end{quote}

is a whole question, built from one new word and two you already own. This is the sentence a market runs on, and you can now say it.

You have also seen \hi{दाम} before, standing in a row of look-alikes beside \textbf{\hi{नाम}} and \textbf{\hi{काम}} when you were learning to pick a name out of a line. Now it means something.

\begin{grammarlens}[title={What you've built}]
Three: \hi{पैसा}, \hi{रुपया}, \hi{दाम} — and a question you can ask with the third.
\end{grammarlens}

\subsection*{Your turn}
Say these aloud:

\begin{itemize}
\raggedright
  \item \emph{dām}
  \item \emph{dām kyā hai?}
  \item it again, pointing at something in the room
\end{itemize}

\begin{tcolorbox}[breakable,colback=teal!4,colframe=teal!35!black,title={Before you move on}]
How do you ask what something costs? (\emph{dām kyā hai?}) Say it once more.
\end{tcolorbox}
\section[mahãgā]{\hi{महँगा} (mahãgā) — expensive}

\label{lesson:HI-C72-expensive}



\begin{quote}
Write \textbf{\hi{दाम}} from memory: \hi{द}, then \hi{ा}, then \hi{म}. Then say the question you built with it.
\end{quote}

\subsection*{You'll want to know: \hi{महँगा}}
\textbf{\hi{महँगा}} (\emph{mahãgā}) — \textquotedblleft{}expensive\textquotedblright{}.

Sanskrit \hi{महार्घ} (\emph{mahārgha}): \textbf{\hi{महा}} is \textquotedblleft{}great\textquotedblright{} and \textbf{\hi{अर्घ}} is \textquotedblleft{}value, price\textquotedblright{}. So the word is \textquotedblleft{}great-priced\textquotedblright{}, which is what it still means.

The moon-and-dot \hi{ँ} over the \hi{ह} sends the vowel through the nose. Say it aloud with the nasal and without it, and you will hear that the nasal is not decoration — it is part of the word.

\hi{महँगा} behaves like \textbf{\hi{बड़ा}} and \textbf{\hi{पुराना}}, which you already own: it comes before the thing it describes and changes its ending to agree with it. You do not have to work the agreement yet. Hear it, and notice the \textbf{-\hi{आ}} ending it shares with those two.

\textbf{\hi{बहुत} \hi{महँगा}} puts last chapter's word to work on this one.

\begin{grammarlens}[title={What you've built}]
Four: a price, and a judgement to pass on it.
\end{grammarlens}

\subsection*{Your turn}
\begin{itemize}
\raggedright
  \item \emph{Say it:} \emph{mahãgā}, with the nose in the middle
  \item \emph{Say it:} \emph{bahut mahãgā}
  \item \emph{Write it:} \hi{महँगा}, putting the \hi{ँ} above the \hi{ह}
\end{itemize}

\begin{tcolorbox}[breakable,colback=teal!4,colframe=teal!35!black,title={Before you move on}]
What does \hi{महँगा} mean? (\textquotedblleft{}expensive\textquotedblright{}.) Say that the price is very high.
\end{tcolorbox}
\section[kharīdnā]{\hi{ख़रीदना} (kharīdnā) — to buy}

\label{lesson:HI-C72-buy}



\begin{quote}
Name the four words of this chapter so far.
\end{quote}

\subsection*{You'll want to know: \hi{ख़रीदना}}
\textbf{\hi{ख़रीदना}} (\emph{kharīdnā}) — \textquotedblleft{}to buy\textquotedblright{}.

Persian \ar{خریدن} (\emph{kharīdan}), \textquotedblleft{}to buy\textquotedblright{}, with the Persian ending swapped for the Hindi \textbf{-\hi{ना}} you now recognise on sight. The whole Persian half of the shopping vocabulary came in this way, over centuries of trade.

The dot under the \hi{ख़} is the nuqta, and it matters. Plain \textbf{\hi{ख}} is the aspirated \emph{k} you met in \textbf{\hi{खेल}}; \textbf{\hi{ख़}} with the dot is a scraped sound further back in the throat, the one Persian and Arabic brought with them. Many printers leave the dot off and the word still reads, but the two letters are not the same.

Put the chapter together and you can shop: \textbf{\hi{दाम} \hi{क्या} \hi{है}?} to ask, \textbf{\hi{बहुत} \hi{महँगा}} to complain, and \hi{ख़रीदना} for what you came to do.

\begin{grammarlens}[title={What you've built}]
Five: \hi{पैसा}, \hi{रुपया}, \hi{दाम}, \hi{महँगा}, \hi{ख़रीदना}. Enough to stand at a counter and settle what a thing costs.
\end{grammarlens}

\subsection*{Your turn}
\begin{itemize}
\raggedright
  \item \emph{Say it:} \emph{kharīdnā}
  \item \emph{Say it:} \emph{dām kyā hai?}, then \emph{bahut mahãgā}
  \item \emph{Write it:} \hi{महँगा} from memory, then \hi{ख़रीदना} beside the printed model
  \item \emph{Say it:} all five in order
\end{itemize}

\begin{tcolorbox}[breakable,colback=teal!4,colframe=teal!35!black,title={Before you move on}]
What does \hi{ख़रीदना} mean? (\textquotedblleft{}to buy\textquotedblright{}.) Ask what a thing costs.
\end{tcolorbox}
